\section{Impossibility of Better ``Distinct-Walk'' Attacks}\label{app:mw}
We assume $h\in \Func([N],[N])$ here for simplicity. 
Zero-walk attacks, as defined Akshima by et al.~\cite{AkshimaCDW20}, assume that an attacker works as follows: in the first unbounded phase, the adversary finds collisions for roughly $S$ points $a_1,a_2,\ldots,a_S\in [N]$ with respect ot $h$, and in the second online phase, it tries to ``walk'' to any of these salts by repeatedly evaluating $h(\cdot,0)$. 

The attack of Akshima et al.~\cite{AkshimaCDW20} that achieves advantage $\Omega({STB}/{(N\log N)})$ falls into the above framework of zero-walk attacks. Akshima et al.\ showed that every such ``zero-walk'' attack could not be used to refute the STB conjecture. To this end, they formalized the zero-walk attacks and considered the best possible variant: an attack where the preprocessing phase examines the functional graph of $h(\cdot,0)$ and select roughly $S$ points $a_1,a_2,\ldots,a_S\in [S]$ such that the number points for which these $S$ points are reachable within $\approx B$ evaluations of $h(\cdot,0)$ is maximized. They showed that the advantage of the best possible zero-walk attack is upper bounded by $O({STB\ln(NB)/N})$.

While this result provides some evidence for the veracity of the STB conjecture, the considered class of attacks ignores the adversary's ability to evaluate $h$ on second inputs other than $0$ which makes it rather limited. In this section, we consider a class of attacks that takes this ability of the adversary into account and consider the other extreme -- at every step $h$ is evaluated on a different message. Here the adversary tries to find a layered tree in the functional graph of $h$ where every layer of nodes is mapped by a distinct $h_i=h(\cdot,i)$ to the next layer. 
We call these attacks the distinct-walk attack: here the preprocessing phase of the attack outputs $S$ salts $a_1,\ldots,a_S$ collisions for these salts \emph{and} $B-2$ \emph{distinct} indices $k_1,\ldots,k_{B-2}$ in $[N]$. The online phase on input $a$ first evaluates $h(a,\hat{\alpha_i})$ for some random $\hat{\alpha_i}$ different from $k_1,\ldots,k_{B-2}$ and then evaluates to get $c_1$. It computes $c_{j}$ by evaluating $h(c_{j-1},k_{j-1})$ for $j\in \{2,\ldots,B-1\}$ and checks if $c_{B-1}\in \{a_1,\ldots,a_S\}$ and returns the appropriate collision if the check succeeds. The reason one might think that this attack might be more effective is the freedom the adversary has in choosing the $B-2$ functions out of $N$. We will however show that this freedom still does not allow the adversary to refute the STB conjecture. This is more evidence in support of the conjecture being true.

We formalize the distinct-walk attack below.
\begin{definition}[$(S,T)$ distinct-walk AI adversary] For $B>2$, a $(S,T)$ distinct-walk AI adversary $\advA=(\advA_1,\advA_2)$ has the following form.
\begin{enumerate}
    \item The first stage adversary $\advA_1$ always produces a bit-encoded output of the form \[\sigma=\{(a_1,\alpha_1,\alpha_1')\ldots,(a_s,\alpha_1,\alpha_s')\},(k_1,k_2,\ldots,k_{B-2})\] where $s=(S-(B-1)\log N)/3\log N$, $a_i\in [N],\alpha_i\alpha_i'\in [N]$, $k_1,k_2,\ldots,k_{B-2}\in [N]$, and $|k_1,k_2,\ldots,k_{B-2}|=B-2$.
    \item The second stage adversary $\advA_2$, on input $a,\sigma$ where $$\sigma= (a_1,\alpha_1,\alpha_1')\ldots,(a_s,\alpha_1,\alpha_s')\}, (k_1,k_2,\ldots,k_{B-2})\;,$$ and given oracle access to $h$ does the following.
    If $\exists i\in [s]:a_i=a$, then output $(\alpha_i,\alpha_i')$. Otherwise for $i=1,\ldots,\lfloor T/B\rfloor$:
    \begin{enumerate}
        \item Choose $\hat{\alpha_i}\getsr [N]\setminus \{k_1,k_2,\ldots,k_{B-2}\}\}$
        \item $c_1\gets h(a,\hat{\alpha_i})$
        \item  $c_j\gets h(c_{j-1},k_{j-1})$ for $j\in 2,\ldots, B-1$:
        \item If $\exists l\in [s]: c_{B-1}= a_l$, then output $$(\hat{\alpha_i}\concat k_{1}\concat \ldots\concat k_{B-2}\concat \alpha_{l},\hat{\alpha_i}\concat k_{1}\concat \ldots\concat k_{B-2}\concat \alpha_{l}').$$
    \end{enumerate}
\end{enumerate}
\end{definition}

In the following theorem we argue that distinct-walk AI adversaries have advantage bounded by $O({STB \ln(NB)}/{N})$, i.e., even these adversaries can not be used to refute the STB conjecture.  The proof of the theorem is largely inspired by the proof of Theorem 8 in~\cite{AkshimaCDW20} which ruled out zero-walk AI attacks refuting the STB conjecture. We give a proof here for the sake of completeness.


\begin{theorem}\label{thm:mfa}
    Let $S,T,B\in \mathbb{N}_{>0}$ such that $SB\geq T$, $B\geq 3$. For any distinct-walk $(S,T)$-AI adversary $\advA$,
    \begin{align*}
    \Adv_{\MD,N,N,B}^{\aicr}(S,T) \leq \frac{500STB\ln(NB)}{N}+\frac{S+1}{N}    \end{align*}
\end{theorem}
\begin{proof}
For $h\in \Func([N],[N])$ define $h_i(x):=h(x,i)$ for $i\in [N]$. We define the ancestors of a salt $a$ with respect to functions $(h_{i_1},\ldots,h_{i_t})$ as follows.
\[
\ancs(a,(h_{i_1},\ldots,h_{i_t})):=\{x: h_{i_t}(h_{i_{t-1}}(\ldots h_{i_1}(x)))=a \}\;.
\] 
We next define for some $t$ the maximum number of $t$-level ancestors of a salt $a$ as follows.
\[
\maxancs(a,h,t):=\max_{{i_1},\ldots,{i_B}\colon |{i_1},\ldots,{i_B}|=B} \ancs(a,(h_{i_1},\ldots,h_{i_t}))\;.
\]
The following lemma is the key lemma we use in the proof of this theorem.
\begin{lemma}\label{lem:keymfa}
    Let $t\in \mathbb{N}_{>0}$ and $h\in \Func([N],[N])$. Then
    \[
    \Prob{\exists a\in [N]:\maxancs(a,h,t)>500t^2\ln (Nt)}\leq 1/N\;.   
    \]
\end{lemma}
First we show how the theorem follows from this lemma.

We say $h$ is bad if $\exists a\in [N]$ such that $\maxancs(a,h,B-2)>500(B-2)^2\ln (N(B-2))$. Otherwise $h$ is good.
%
Observe that the online phase can only succeed if either $a\in \{a_1,\ldots,a_s\}$ which happens with probability at most $S/N$ or in any of the $T/B$ iterations $c_1$ is in $\ancs(a_i,(h_{k_1},\ldots,h_{k_{B-2}}))$ for some $i\in [s]$. Given that $h$ is good, we have that
\begin{align*}
|\bigcup_{i=1}^s \ancs(a_i,(h_{k_1},\ldots,h_{k_{B-2}}))| & \leq 500S(B-2)^2\ln (N(B-2)) \\ & \leq 500SB^2\ln (NB).
\end{align*}
Using total probability, we have 
\begin{align*}
\Adv_{\MD,N,N,B}^{\aicr}(S,T) &\leq \frac{S}{N}+\Pr[\advA \text{ wins}\mid h \text{ is good}]+\Pr[h\text{ is bad}]\\
&\leq \frac{500STB\ln(NB)}{N}+\frac{S+1}{N}\;.
\end{align*}
\end{proof}
We next prove the key \Cref{lem:keymfa}.
\begin{proof}[\Cref{lem:keymfa}]
 Fix $k_1,\ldots,k_t\in [N]$ such that $|\{k_1,\ldots,k_t\}|=t$. Fix $a\in N$.
    Let $\Ne^{(0)}(a):=\{a\}$ and for $i=1,\ldots,t$
    \[
    \Ne^{(i)}(a):=\{x: h_{k_i}(y)=x\text{ for }y\in \Ne^{(i)}(a)  \}\;.    
    \]
    \begin{myclaim}\label{claim:mw}
    Let $t\in \mathbb{N}_{>0}$. For all $a\in [N]$ the following holds.
    \[
    \Pr\left[|\Ne^{(t)}(a)|> 500t^2 \ln(Nt)\right]\leq N^{-3t} \;.    
    \]
    \end{myclaim}
    First it is easy to see that if the claim is true, taking an union bound over all choices of $a$ and all choices of $k_1,\ldots,k_t\in [N]$ we have that
    \begin{align*}
        \Prob{\exists a\in [N]:\maxancs(a,h,t)>500t^2\ln (Nt)}&\leq  N\cdot \binom{N}{t}t!N^{-3t}\\
        &\leq  N\cdot N^t\cdot N^{-3t}\leq N^{-t}\leq 1/N\;.   
    \end{align*}
\end{proof}
    So we next prove Claim~\labelcref{claim:mw}.
    \begin{proof} [Claim~\labelcref{claim:mw}]
    Let $P_j:=|\Ne^{{j}}(a)|$. The following two claims together imply claim~\labelcref{claim:mw}.
    \begin{myclaim}\label{claim:one}
        The probability that there exists $j$ such that $P_{j}<50 t^2\ln (Nt)$ and $P_{j+1}>100 t^2\ln (Nt)$ is at most $(Nt)^{-16t^2}$.  
    \end{myclaim}
    \begin{myclaim}\label{claim:two}
        The probability that there exist $k,l$ ($k<l$) such that $50t^2\ln (Nt)<P_k<100 t^2\ln (Nt)$, for all $j\in \{k+1,\ldots,l-1\}$ $P_{j}>50t^2\ln (Nt)$  and $P_{l}>500 t^2\ln (Nt)$ is at most $N^{-4t}$.
    \end{myclaim}
    Combining Claims~\labelcref{claim:one,claim:two}, we have that
    \begin{align*}
    \Pr[\Ne^{(t)}(a)>500t^2 \ln (Nt)]& =    \Pr[P_t>500t^2 \ln (Nt)]\\& \leq N^{-4t}+(Nt)^{-16t^2}\leq N^{-3t}\;.
    \end{align*}
    This proves Claim~\labelcref{claim:mw}.
\end{proof}
We next prove the two claims~\labelcref{claim:one,claim:two}.
\begin{proof}[Claim~\labelcref{claim:one}]  Consider the indicator random variables $X_a$ defined as follows.
    \[
    X_a=\begin{cases}
        1 & \text{ if }a\in \Ne^{(j)}(a)\\
        0 &\text{ otherwise}
    \end{cases}    
    \]
    Let $X=\sum_{a\in [N]}X_a$. It is easy to see that all the $X_a$'s are independent, $X=P_j$ and $\Exp{X}=P_{j-1}$. Using the multiplicative Chernoff bound (Proposition~\labelcref{prop:chernmult}) we have
    \[
    \Pr[X>(1+\delta)\Exp{X}] \leq \exp\left(\frac{-\delta^2\Exp{X}}{2+\delta}\right)\;.
    \]
    Here $\delta\Exp{X}\geq 50t^2\ln(Nt)$. Also we have that
    $\delta\geq 1$.
    This means
    $\delta/(2+\delta)=1-2/(2+\delta)\geq 1/3$.
    Hence
    \[
        \Pr[X>(1+\delta)\Exp{X}] \leq \exp\left(\frac{-50t^2\ln(Nt)}{3}\right)\leq \frac{1}{(Nt)^{16t^2}}\;.
    \]
\end{proof}
    % If there does not exist $P_j$ $P_t>500t^2 \ln(Nt)$ with very high probability there must exist some $j,j'$ such that $50t^2\ln (Nt)<P_j<100 t^2\ln (Nt)$ 
\begin{proof}[Claim~\labelcref{claim:two}]
    Define random variables $\delta_j$ as such that 
    \[
    P_{j}=(1+\delta_j)P_{j-1}\;.    
    \]
    For $j\in \{l+1,\ldots,k\}$ using the multiplicative Chernoff bound (Proposition~\labelcref{prop:chernmult}) we get
    \begin{align}\label{eq:delta}
    \Pr[\delta_j\geq \delta]& =\Pr[P_{j}\geq (1+\delta)P_{j-1}]\nonumber\\&\leq \exp\left(\frac{-\delta^2P_{j-1}}{2+\delta}\right)\leq \exp\left(\frac{-\delta^250t^2\ln (Nt)}{2+\delta}\right)\;.
    \end{align}
    From the premise of the claim, it follows that $\prod_{j=l+1}^k(1+\delta_j)\geq 5$. Using $1+x\geq e^x$ we have that 
    \[
    e^{\sum_{j=l+1}^k\delta_j}\geq 5\;.
    \]
    It follows that 
    $\sum_{j=l+1}^k\delta_j\geq \ln 5$.
    We want to next compute $\Pr[\sum_{j=l+1}^k\delta_j\geq \ln 5]$. Observe that this is same as the following probability.
    \[
        \Pr[N^{3t\sum_{j=l+1}^k\delta_j}\geq N^{3t\ln 5}]\;.
    \]
    We will prove that for each $j\in \{l+1,\ldots,k\}$ 
    \begin{equation}\label{eq:ineq}
        \Pr[N^{3t\delta_j}\geq N^{3t\ln 5}]\leq N^{2-3t\ln 5}\;.
    \end{equation}
    Using the union bound and $t\leq N$, it follows that 
    \[
        \Pr[N^{3t\sum_{j=l+1}^k\delta_j}\geq N^{3t\ln 5}]\leq N^{3-3t\ln 5}\leq N^{-2t\ln 5}\leq N^{-4t}\;.
    \]
    So, we need to now show that the inequality in \Cref{eq:ineq} holds.
    Using Markov's inequality (Proposition ~\labelcref{prop:markov}), we have that 
    \begin{equation}\label{eq:second}
        \Pr[N^{3t\delta_j}\geq N^{3t\ln 5}]\leq \frac{\Exp{N^{3t\delta_j}}}{N^{3t\ln 5}}\;.
    \end{equation}
    From the definition of expectation we have
    \begin{align}
        \Exp{N^{3t\delta_j}}&\leq \int_0^\infty \Pr[N^{3t\delta_j}\geq x]~dx \nonumber\\
        &= \int_0^{N} \Pr[N^{3t\delta_j}\geq x]~dx+\int_N^{N^{3t}} \Pr[N^{3t\delta_j}\geq x]~dx +\int_{N^{3t}}^\infty \Pr[N^{3t\delta_j}\geq x]~dx \nonumber\\
        &\leq N+\int_N^{N^{3t}} \Pr[N^{3t\delta_j}\geq x]~dx +\int_{N^{3t}}^\infty \Pr[N^{3t\delta_j}\geq x]~dx\;.\label{eq:exp}
    \end{align}
    The second step follows from the splitting the sum and the third step follows since $\Pr[N^{3t\delta_j}\geq x]\leq 1$.
    Note that $\Pr[N^{3t\delta_j}\geq x$ is same as $\Pr[\delta_j\geq \ln x/(3t\ln N)]$. Let $\ln x/(3t\ln N)=\delta$.
    
    We first upper bound $\int_N^{N^{3t}} \Pr[N^{3t\delta_j}\geq x]~dx$. In the range of $x$ for this integral, $\delta=\ln x/(3t\ln N) \leq 1$. Hence, using \Cref{eq:delta}
    \begin{align*}
        \Pr[\delta_j\geq \delta]\leq \exp\left(\frac{-\delta^2\cdot 50t^2\ln (Nt)}{2+\delta}\right)
        &\leq \exp\left(\frac{-(\ln x)^2\cdot 50t^2\ln (Nt)}{3\cdot 9t^2(\ln N)^2}\right)\\
        &\leq \exp(-\frac{50\ln x}{27})\leq x^{-50/27}\;.
    \end{align*}
    The third inequality follows since $\ln x\geq \ln N$ and $\ln (Nt)\geq \ln N$.
%
This implies that 
 \begin{equation}\label{eq:intf}
        \int_N^{N^{3t}} \Pr[N^{3t\delta_j}\geq x]~dx\leq  \int_N^{N^{3t}} x^{-50/27}~dx \leq \frac{27N^{-23/27}}{23}\;.
    \end{equation}
    We next upper bound $\int_{N^{3t}}^\infty \Pr[N^{3t\delta_j}\geq x]~dx$. In the range of $x$ for this integral, $\delta \geq 1$. Hence, using \Cref{eq:delta}
    \begin{align*}
        \Pr[\delta_j\geq \delta]\leq \exp\left(\frac{-\delta^2 \cdot 50t^2\ln (Nt)}{2+\delta}\right)&\leq \exp\left(\frac{-\delta \cdot 50t^2\ln (Nt)}{3}\right)\\ & =\exp\left(\frac{-\ln x \cdot 50t^2\ln (Nt)}{3t\ln N}\right)\leq x^{-16t}\;.
    \end{align*}
    The second inequality follows since $\delta\geq 1$ and the third inequality follows because $\ln (Nt)\geq \ln N$.

   This implies that
    \begin{equation}\label{eq:ints}
        \int_{N^{3t}}^\infty \Pr[N^{3t\delta_j}\geq x]~dx\leq  \int_{N^{3t}}^\infty x^{-16t}~dx \leq  N^{-t}\;.
    \end{equation}
    From~\Cref{eq:exp,eq:intf,eq:ints} it follows 
    \[\Exp{N^{3t\delta_j}}\leq N+\frac{27N^{-23/27}}{23}+N^{-t}\leq N^2\;.\]
Finally, from~\Cref{eq:second} we can conclude that
\[\Pr[N^{3t\delta_j}\geq N^{3t\ln 5}] \leq N^{2-3t\ln 5}\;.\]
\end{proof}


